% Volume V: Admissibility and Continuation
% Part X. Public Institutions as Epistemic Ecologies
% Chapter 57: Bureaucracy and Slow Knowledge
% Spine: proof-spines/ch057-spine.md

\chapter{Bureaucracy and Slow Knowledge}
\label{ch:ch057}

Bureaucratic delay can conceal indifference or protect against haste. This chapter distinguishes \textbf{pathological delay}, \textbf{capacity delay}, \textbf{verification delay}, \textbf{jurisdictional delay}, and \textbf{strategically imposed delay} instead of treating slowness as a single phenomenon.

A normative theory of \textbf{slow knowledge} must ask what work occurs during the interval, who bears its costs, and whether the process remains inspectable. Delay is justified when it performs comparison, checking, routing, or error control that cannot be skipped without loss; it becomes suspect when time is consumed without epistemic work or when the costs of waiting are invisibly shifted onto those already exposed to institutional power.
